\documentclass[preprint,12pt]{elsarticle}

% Basic packages
\usepackage[a4paper, margin=1in]{geometry}
\usepackage{graphicx}
\usepackage{amsmath}
\usepackage{hyperref}
\usepackage{longtable}
\usepackage{booktabs}
\usepackage{caption}
\usepackage{subcaption}
\captionsetup[table]{skip=10pt}

% Add Unicode support
\usepackage[utf8]{inputenc}
\usepackage[T1]{fontenc}
\usepackage{textcomp}

% Add color support
\usepackage{xcolor}

% Line numbering
\usepackage{lineno}
\leftlinenumbers
\renewcommand\linenumberfont{\normalfont\small\color{gray}}

% Add support for text formatting
\usepackage{array}
\usepackage{enumitem}
\usepackage{float}
\usepackage{gensymb}

% Supplementary numbering
\renewcommand{\thesection}{S\arabic{section}}
\renewcommand{\thetable}{S\arabic{table}}
\renewcommand{\thefigure}{S\arabic{figure}}
\renewcommand{\theequation}{S\arabic{equation}}

\biboptions{authoryear,round,semicolon}

%-----------------------------------------------
% METADATA
\title{Supplementary Material for: Predicting plant diversity facets from Landsat phenology and topography in central and southern Chile}

\author[1,2,3]{Javier Lopatin\corref{cor1}}
\ead{javier.lopatin@uai.cl}

\author[4]{Rocío Araya-López}
\ead{ro.arayalopez@gmail.com}

\author[2,5]{Dylan Craven}
\ead{dylan.craven@umayor.cl}

\address[1]{Center for Earth and Space, Faculty of Engineering and Science, Adolfo Ibáñez University, Av. Diagonal Las Torres 2460, Santiago, Chile, 7941169}

\address[2]{Data Observatory Foundation, ANID Technology Center No. DO210001, Santiago, Chile, 7510277}

\address[3]{Center for Climate Resilience Research (CR)2, University of Chile, Santiago, Chile, 8370449}

\address[4]{Deakin  Marine Research and Innovation Centre, School of Life and Environmental Sciences, Deakin University, Burwood Campus, Burwood, VIC 3125, Australia}

\address[5]{GEMA Center for Genomics, Ecology and the Environment, Universidad Mayor, Santiago, Chile}

\cortext[cor1]{Corresponding author; javier.lopatin@uai.cl}

%-----------------------------------------------
\begin{document}
\begin{frontmatter}
\end{frontmatter}

\linenumbers

\section{Choice of the vegetation index}
\label{sec:s_index}

The main text uses the kernel NDVI (kNDVI) as the phenological substrate. We selected it
among five vegetation indices computed from the same Landsat surface-reflectance
observations: the normalised difference vegetation index \citep{tucker1979red}, the
enhanced vegetation index \citep{huete2002overview}, kNDVI \citep{campsvalls2021kndvi},
the normalised burn ratio \citep{key2006landscape} and the soil-adjusted vegetation index
\citep{huete1988savi}:
\begin{align}
\mathrm{NDVI} &= \frac{\rho_\mathrm{NIR} - \rho_\mathrm{red}}{\rho_\mathrm{NIR} + \rho_\mathrm{red}}, &
\mathrm{kNDVI} &= \tanh\!\left(\mathrm{NDVI}^2\right), \\
\mathrm{NBR} &= \frac{\rho_\mathrm{NIR} - \rho_\mathrm{SWIR2}}{\rho_\mathrm{NIR} + \rho_\mathrm{SWIR2}}, &
\mathrm{SAVI} &= \frac{1.5\,(\rho_\mathrm{NIR} - \rho_\mathrm{red})}{\rho_\mathrm{NIR} + \rho_\mathrm{red} + 0.5}, \\
\mathrm{EVI} &= \frac{2.5\,(\rho_\mathrm{NIR} - \rho_\mathrm{red})}{\rho_\mathrm{NIR} + 6\rho_\mathrm{red} - 7.5\rho_\mathrm{blue} + 1}.
\end{align}

For each index we built the raw three-year series of the plot pixel, folded it into the
serpentine image and trained the two-dimensional network under the same context block,
targets, spatial block cross-validation and five seeds as in the main text (Methods,
Sections~1.3 to~1.8). We summarised each index by the mean out-of-fold $R^2$ over the three
facets with usable signal (LCBD, PD$_0$ and TD$_0$), computing that mean within each seed
first, so that its standard deviation is the seed-to-seed dispersion of the summary itself.

The five indices are indistinguishable (Table~\ref{tab:s_index}). Their summary means span
0.585 to 0.605, and the largest gap between any index and kNDVI (0.016, for NBR) is smaller
than the two standard deviations combined in quadrature (0.017); the gap between the three
leading indices is two orders of magnitude smaller than their dispersion. The same was true
of the earlier configuration of this sweep, where the leading margin was 0.0007 against a
seed standard deviation of 0.016. We therefore do not rank the indices, and we note that
with five seeds the standard deviations themselves are uncertain by roughly a third.

We retained kNDVI for two reasons that do not depend on that ranking. It is the index of
the masked-autoencoder checkpoint used to initialise the deployed network, so keeping it
avoids a mismatch between pretraining and fine-tuning; and on TD$_0$, the facet with the
strongest signal, kNDVI has the highest mean (0.783) and the narrowest between-seed spread
(0.014, against 0.029 to 0.066 for the other four), although that lead is itself within
dispersion. The between-fold spread does not separate the indices: computed per fold, every
index scores about 0.4 on TD$_0$ with a standard deviation of the same order as the mean,
because the 20-km blocks differ strongly among themselves.
No index produced consistent positive scores on the abundance-weighted facets.

\begin{table}[t]
\centering
\footnotesize
\caption{Out-of-fold $R^2$ by vegetation index (two-dimensional network, raw series of the
plot pixel, topography and plot context, spatial block cross-validation). Mean $\pm$
standard deviation over five seeds. The last column is the mean over LCBD, PD$_0$ and
TD$_0$, taken within each seed before averaging. Indices are ordered by that summary; the
differences between them are smaller than their dispersion.}
\label{tab:s_index}
\begin{tabular}{lcccccc}
\toprule
Index & LCBD & PD$_0$ & PD$_1$ & PD$_2$ & TD$_0$ & Mean \\
\midrule
SAVI  & $0.427\pm0.004$ & $0.616\pm0.006$ & $-0.021\pm0.038$ & $0.025\pm0.028$ & $0.772\pm0.029$ & $0.605\pm0.009$ \\
NDVI  & $0.437\pm0.007$ & $0.604\pm0.034$ & $0.014\pm0.030$ & $0.065\pm0.016$ & $0.761\pm0.060$ & $0.601\pm0.020$ \\
kNDVI & $0.433\pm0.004$ & $0.586\pm0.023$ & $0.010\pm0.029$ & $0.047\pm0.023$ & $\mathbf{0.783\pm0.014}$ & $0.601\pm0.007$ \\
EVI   & $0.424\pm0.009$ & $0.605\pm0.016$ & $-0.006\pm0.036$ & $0.036\pm0.032$ & $0.753\pm0.066$ & $0.594\pm0.021$ \\
NBR   & $0.440\pm0.010$ & $0.599\pm0.019$ & $0.018\pm0.033$ & $0.071\pm0.023$ & $0.716\pm0.051$ & $0.585\pm0.016$ \\
\bottomrule
\end{tabular}

\vspace{2pt}
\footnotesize TD$_1$ and TD$_2$ are omitted from the table because every index scores
between $-0.17$ and $-0.06$ on them; the full values are in the repository.
\end{table}

% Bibliography settings
\bibliographystyle{elsarticle-harv}
\bibliography{refs}

\end{document}
